\section{Návrh a~architektúra systému}
\label{sec:architektura}

\emph{Kapitola obsahuje prvý návrh architektúry. Jednotlivé časti budú
spresnené podľa výsledkov experimentov.}

Jadrom systému je \textbf{orchestrátor}, komponent, ktorý na plánovanie
využíva hlavný jazykový model. Audit AI modelu rozloží na menšie podúlohy
a~prideľuje ich \textbf{agentom}; agenti na svoje úlohy využívajú menšie
jazykové modely s~malým kontextom. \emph{Statickí} agenti sú potrební pri
každom audite a~bežia stále, \emph{dynamickí} sa spúšťajú a~zastavujú
podľa potreby. Systém nie je viazaný na konkrétne prostredie: beží na
ľubovoľnom počítači s~GPU a~kontajnerovým runtime, napríklad na serveri
školy alebo firmy.

%%%%%%%%%%%%%%%%%%%%%%%%%%%%%%%%%%%%%%%%
\subsection{Princípy návrhu}
\label{sec:principy}

\begin{itemize}
    \item \textbf{Všetko lokálne.} Systém nepoužíva žiadne externé API.
    Jazykové modely sú stiahnuté na hostiteľskom počítači a~bežia
    v~\texttt{llama-server} z~projektu llama.cpp~\cite{llamacpp}. Testovaný
    model sa nahrá ako súbory.
    \item \textbf{Internet iba výnimočne.} Prístup von je vypnutý; ak je
    potrebný, ide cez proxy na niekoľko overených stránok (whitelist,
    tabuľka~\ref{tab:whitelist}).
    \item \textbf{Žiadne zistenie bez dôkazu.} Každé tvrdenie v~profile
    odkazuje na konkrétny úryvok zdroja, inak sa údaj označí ako chýbajúci.
    \item \textbf{Malé úlohy namiesto jednej veľkej.} Agent rieši úzku
    podúlohu s~malým kontextom, čo znižuje riziko halucinácií.
    \item \textbf{Kontrola kritikom.} Každé tvrdenie, ktoré vytvorí
    jazykový model, overí kritik (vzor actor-critic
    a~self-refine~\cite{madaan2023selfrefine}).
\end{itemize}

\paragraph{Whitelist.} Zberač má cez proxy prístup iba na stránky
v~tabuľke~\ref{tab:whitelist}. Pri Hugging Face a~GitHub sa uprednostňujú
oficiálne repozitáre vývojárov modelu; informácie z~ostatných repozitárov
treba dodatočne overiť.

\begin{table}[ht]
    \centering
    \caption{Whitelist overených zdrojov pre zberača}
    \label{tab:whitelist}
    \small
    \begin{tabular}{@{}>{\raggedright\arraybackslash}p{2.4cm}>{\raggedright\arraybackslash}p{3.2cm}>{\raggedright\arraybackslash}p{8.6cm}@{}}
        \toprule
        \textbf{Zdroj} & \textbf{Doména} & \textbf{Účel} \\
        \midrule
        Hugging Face & \texttt{huggingface.co} & model cards, technické
            parametre, konfigurácie a~informácie o~modeloch \\
        GitHub & \texttt{github.com} & oficiálne repozitáre, zdrojový kód
            a~technická dokumentácia \\
        arXiv & \texttt{arxiv.org} & vedecké publikácie, výskumné články
            a~výsledky testovania modelov \\
        NIST & \texttt{nist.gov} & bezpečnosť AI, hodnotenie rizík
            a~odporúčania pre dôveryhodné AI systémy \\
        \bottomrule
    \end{tabular}
\end{table}

\paragraph{Skill.} Skill je opis jednej schopnosti agenta, ktorý tvoria
tri časti: inštrukcia pre jazykový model (prompt), zoznam povolených
nástrojov a~JSON schéma výstupu. Skilly sú uložené v~registri skillov;
orchestrátor ku každej podúlohe vyberie vhodný skill a~odovzdá ho agentovi
spolu s~podúlohou. V~experimentálnom variante môže generátor skillov
vytvoriť nový skill, ak žiadny z~existujúcich nevyhovuje; systém sa
porovná s~touto možnosťou aj bez nej.

\paragraph{Prečo kontajnery.} Samostatný kontajner je nutný najmä pre
analytika artefaktov, ktorý číta súbory testovaného modelu: súbory vo
formáte pickle môžu pri načítaní spustiť ľubovoľný kód~\cite{hf-pickle},
preto sa spracúvajú izolovane a~bez prístupu na internet (komunikuje iba
s~databázou vo vnútornej sieti kontajnerov). Zberač ako jediný potrebuje
internet, preto beží oddelene za proxy. Rovnako izolovane, bez prístupu
na internet, beží tester, ktorý testovaný model spúšťa. Agenti, ktorí iba volajú
jazykový model, sú v~kontajneroch kvôli jednotnému nasadeniu a~nezávislému
spúšťaniu; ak sa to ukáže ako zbytočná réžia, môžu bežať ako procesy
v~kontajneri orchestrátora.

%%%%%%%%%%%%%%%%%%%%%%%%%%%%%%%%%%%%%%%%
\subsection{Komponenty systému}
\label{sec:komponenty}

Obrázok~\ref{fig:komponenty} zachytáva štruktúru systému ako UML diagram
komponentov. Plné šípky s~popisom nie sú závislosti v~zmysle UML, ale
ukazujú \emph{tok dát} medzi komponentmi. Značka \textsf{AI} označuje
komponenty, ktoré volajú jazykový model. Úložisko dôkazov je databáza,
testované modely a~katalóg kritérií sú dáta (artefakty). Audit sa zatiaľ
spúšťa z~príkazového riadka (CLI). Prerušovanou čiarou je vyznačená
experimentálna časť s~generátorom skillov.

\begin{figure}[tp]
\centering
\includestandalone[width=\textwidth]{uml/komponenty}
\caption{UML diagram komponentov systému (šípky znázorňujú tok dát)}
\label{fig:komponenty}
\end{figure}

%%%%%%%%%%%%%%%%%%%%%%%%%%%%%%%%%%%%%%%%
\subsection{Orchestrátor a~agenti}
\label{sec:agenti}

Obrázok~\ref{fig:agenti} je prehľadový diagram (nie UML), ktorý pre
orchestrátor a~každého agenta uvádza vstup, úlohu jazykového modelu
a~výstup s~jeho cieľom. Výstupy agentov s~jazykovým modelom idú ku
kritikovi. Výstupy bez jazykového modelu, teda nálezy deterministických
kontrol analytika artefaktov a~dokumenty stiahnuté zberačom, sa ukladajú
priamo do úložiska; otázky tvorcu otázok idú cez orchestrátor
používateľovi.

Pre dynamické overenie atribútov (kapitola~\ref{sec:specifikacia-atributov})
pribudne dynamický agent \emph{tester}. Spustí testovaný model
v~izolovanom kontajneri a~vykoná na ňom testy nástrojmi ako garak, PyRIT
alebo Inspect AI (tabuľka~\ref{tab:nastroje}); výsledky testov ukladá do
úložiska ako dôkazy.

\begin{figure}[ht]
\centering
\includestandalone[width=\textwidth]{uml/agenti}
\caption{Prehľadový diagram orchestrátora a~agentov}
\label{fig:agenti}
\end{figure}

%%%%%%%%%%%%%%%%%%%%%%%%%%%%%%%%%%%%%%%%
\subsection{Informačný model profilu}
\label{sec:profil}

Obrázok~\ref{fig:profil} zachytáva jednotný informačný model profilu ako
UML diagram tried. Základnou jednotkou nie je hodnota, ale
\emph{tvrdenie}: pole profilu s~hodnotou, stavom (uvedené, odvodené,
chýbajúce, rozporné) a~mierou istoty. Atribút \emph{oblasť} priraďuje
tvrdenie k~jednej z~oblastí profilu podľa zadania (model card, účel
a~obmedzenia, dáta, evaluácie, bezpečnosť). Tvrdenie so stavom uvedené
alebo odvodené musí mať aspoň jeden dôkaz, teda úryvok s~polohou
v~konkrétnom zdroji; toto obmedzenie je na diagrame zapísané pri triede
Tvrdenie. Jeden úryvok môže dokladať viac tvrdení. Riziká, hodnotenia
požiadaviek a~doplňujúce otázky sú súčasťou profilu; odpoveď používateľa
na otázku sa uloží ako nový zdroj. Zoznam polí profilu (atribúty
extrahované z~model cards) bude doplnený.

\begin{figure}[ht]
\centering
\includestandalone[width=\textwidth]{uml/profil}
\caption{UML diagram tried informačného modelu profilu}
\label{fig:profil}
\end{figure}

%%%%%%%%%%%%%%%%%%%%%%%%%%%%%%%%%%%%%%%%
\subsection{Nasadenie}
\label{sec:nasadenie}

Obrázok~\ref{fig:nasadenie} ukazuje rozmiestnenie komponentov na
hostiteľskom počítači. Každá časť systému je samostatný kontajner (Docker
alebo Podman). Používateľ ovláda systém cez CLI, lokálne alebo vzdialene;
jazykové modely, databáza aj nahraté testované modely sú dostupné iba
vnútri hostiteľa. Jazykové modely obsluhuje jedna alebo viac inštancií
\texttt{llama-server}, keďže jedna inštancia obsluhuje jeden model
(hlavný model orchestrátora a~menšie modely agentov). Register skillov je
uložený v~databáze, generátor skillov a~generátor reportu bežia
v~kontajneri orchestrátora. Analytik artefaktov nemá prístup na internet; s~databázou komunikuje iba
cez vnútornú sieť kontajnerov.

\begin{figure}[ht]
\centering
\includestandalone[width=\textwidth]{uml/nasadenie}
\caption{UML diagram nasadenia}
\label{fig:nasadenie}
\end{figure}

%%%%%%%%%%%%%%%%%%%%%%%%%%%%%%%%%%%%%%%%
\subsection{Priebeh auditu}
\label{sec:pipeline}

Obrázok~\ref{fig:pipeline} opisuje priebeh auditu jedného modelu ako UML
stavový diagram. Každé tvrdenie AI agenta prechádza kontrolou kritika;
zamietnutý výstup opraví ten agent, ktorý ho vytvoril. Audit má limity:
počet opráv jedného výstupu, počet iterácií a~čas čakania na odpoveď
používateľa. Po ich vyčerpaní sa audit ukončí reportom, v~ktorom zostanú
neoverené tvrdenia a~nezodpovedané otázky označené ako otvorené. Report
obsahuje aj odporúčania, ako zlepšiť model a~jeho dokumentáciu: ktoré
údaje doplniť a~ktoré zistené riziká zmierniť.

\begin{figure}[ht]
\centering
\includestandalone[width=\textwidth]{uml/pipeline}
\caption{UML stavový diagram priebehu auditu}
\label{fig:pipeline}
\end{figure}
